% =====================================================================
%  CARNET D'ENTRAÎNEMENT — Fiche 12 (Chimie organique)
%  THÈME 3 — Nomenclature et priorité des fonctions
%  NIVEAU FACILE — ÉNONCÉS
% =====================================================================
\documentclass[11pt,a4paper]{article}
\usepackage[utf8]{inputenc}
\usepackage[T1]{fontenc}
\usepackage{lmodern}
\usepackage[french]{babel}
\usepackage{amsmath,amssymb,mathtools}
\providecommand{\degree}{^\circ}
\usepackage{icomma}
\usepackage{siunitx}
\sisetup{output-decimal-marker={,},per-mode=symbol}
\usepackage[version=4]{mhchem}
\usepackage{chemfig}
\setchemfig{atom sep=2em,bond length=0.9em,cram width=2.5pt}
\usepackage{xcolor}
\usepackage{array}
\usepackage{booktabs}
\usepackage{tabularx}
\usepackage[most]{tcolorbox}
\usepackage{enumitem}
\usepackage{tikz}
\usetikzlibrary{arrows.meta,positioning,calc}
\usepackage{fancyhdr}
\usepackage{titlesec}
\usepackage[hidelinks]{hyperref}
\usepackage[a4paper,margin=2.1cm,headheight=26pt,headsep=9pt]{geometry}

\definecolor{primary}{RGB}{150,98,162}
\definecolor{primarydark}{RGB}{92,52,110}
\definecolor{excol}{RGB}{0,135,90}
\definecolor{attncol}{RGB}{200,40,55}
\definecolor{methcol}{RGB}{90,90,100}

\newcommand{\runtitle}{Thème 3 — Facile — Énoncés}
\pagestyle{fancy}\fancyhf{}
\fancyhead[L]{\small\textcolor{primarydark}{\textbf{Fiche 12} — Chimie organique}}
\fancyhead[R]{\small\textcolor{primarydark}{\runtitle}}
\fancyfoot[C]{\small\thepage}
\renewcommand{\headrulewidth}{0.8pt}
\renewcommand{\headrule}{\hbox to\headwidth{\color{primary}\leaders\hrule height \headrulewidth\hfill}}
\setlength{\parindent}{0pt}\setlength{\parskip}{3pt}

\tcbset{boxbase/.style={enhanced,breakable,boxrule=0.9pt,arc=2.5pt,
   left=3mm,right=3mm,top=2mm,bottom=2mm,fonttitle=\bfseries,coltitle=white,
   toptitle=0.5mm,bottomtitle=0.5mm}}
\newtcolorbox[auto counter]{exo}[1]{boxbase,colback=primary!4,colframe=primary,
   title={\textsc{Exercice}~\thetcbcounter\ \textmd{--- \itshape #1}}}
\newtcolorbox{consigne}{boxbase,colback=methcol!8,colframe=methcol,sharp corners,
   title={\textsc{Consignes}}}

\newcommand{\banner}[2]{%
\begin{tcolorbox}[enhanced,colback=primary,colframe=primary,arc=5pt,boxrule=0pt,
   left=5mm,right=5mm,top=2.5mm,bottom=2.5mm,coltext=white]
{\large\bfseries #1}\\[1pt]{\small #2}
\end{tcolorbox}\vspace{2pt}}

\begin{document}

\banner{Thème 3 — Nomenclature et priorité des fonctions}{Niveau \textbf{Facile} — Énoncés \quad$\vert$\quad Fiche 12, Chimie organique}

\begin{consigne}
Exercices \textbf{ouverts}. Rappel de l'échelle de priorité : $\ce{-COOH} > \ce{-COOR} > \ce{-CONH2} > \ce{-CHO} > \text{C=O} > \ce{-OH} > \ce{-NH2} > \ce{-O-} > \text{C=C} > \text{C}\equiv\text{C}$.
\end{consigne}

% ---------------------------------------------------------------------
\begin{exo}{donner le nom (fonction unique)}
Nomme chacune des molécules suivantes (une seule fonction à chaque fois) :
\begin{center}\renewcommand{\arraystretch}{1.2}
\begin{tabular}{cccc}
\chemfig{CH_3-[:30]CH_2-[:-30]CH_2-[:30]OH} &
\chemfig{CH_3-[:30]CH_2-[:-30]CH_2-[:30](=[:90]O)-[:-30]H} &
\chemfig{CH_3-[:30](=[:90]O)-[:-30]CH_3} &
\chemfig{CH_3-[:30]CH_2-[:-30](=[:-90]O)-[:30]OH} \\
\textbf{(A)} & \textbf{(B)} & \textbf{(C)} & \textbf{(D)} \\
\end{tabular}
\end{center}
\end{exo}

% ---------------------------------------------------------------------
\begin{exo}{numéroter la chaîne}
On donne l'\textbf{acide butanoïque} :
\begin{center}\chemfig{CH_3-[:30]CH_2-[:-30]CH_2-[:30](=[:90]O)-[:-30]OH}\end{center}
Quel carbone porte le \textbf{n\textsuperscript{o} 1} ? Numérote les 4 carbones de la chaîne et justifie la règle utilisée.
\end{exo}

% ---------------------------------------------------------------------
\begin{exo}{repérer la fonction principale}
La molécule ci-dessous porte \textbf{deux} fonctions :
\begin{center}\chemfig{H-[:30](=[:90]O)-[:-30]CH_2-[:30]CH_2-[:-30]OH}\end{center}
\begin{enumerate}[label=\textbf{\alph*)},leftmargin=2em,itemsep=1pt]
  \item Identifie les deux fonctions.
  \item En utilisant l'échelle de priorité, indique laquelle est la \textbf{fonction principale} (donne le suffixe) et laquelle devient un \textbf{préfixe}.
\end{enumerate}
\end{exo}

% ---------------------------------------------------------------------
\begin{exo}{du suffixe à la fonction}
Pour chacun des noms suivants, indique la \textbf{fonction} révélée par le suffixe :
\begin{center}\small
butan-2-\textbf{ol} \quad$\vert$\quad propan\textbf{al} \quad$\vert$\quad acide éthano\textbf{ïque} \quad$\vert$\quad butan-2-\textbf{one} \quad$\vert$\quad éthan\textbf{amide}
\end{center}
\end{exo}

% ---------------------------------------------------------------------
\begin{exo}{nommer une molécule ramifiée}
Nomme la molécule ci-dessous (repère la fonction, la chaîne principale, la position du groupe $\ce{-OH}$ et de la ramification $\ce{-CH3}$) :
\begin{center}\chemfig{CH_3-[:30]CH(-[:-90]CH_3)-[:-30]CH_2-[:30]OH}\end{center}
\end{exo}

\end{document}
